% Second pass (Claude, 2026-11-10): first pass checked against the archive images; changes in LM_C4-039-057_changes.md.
\scanpage{39}
\mspage{667}

\[
\tag{24} l_1^\nu=\lambda_1^\nu\omega^{\nu_2}.
\]

Setting (cf. (19))
\[
\tag{25}
\left\{\begin{aligned}
\beta&=\beta_0\tau^{\nu_2}\qquad(\unsure{\beta_0|\beta}\in\Gamma\underline M)\\
\eta_0&=\beta_0\sigma(\beta_0)^{-1}
\end{aligned}\right.
\]
we shall thus have
\[
\eta\overset{\unsure{\mathrm{dfn}}}{=}
\beta\sigma(\beta)^{-1}
=\beta_0\sigma(\beta_0)\omega^{\nu_2}
=\eta_0\omega^{\nu_2}
\]
and the loop equation takes the form
of an equation in $\underline M$
\[
\tag{26}
\unsure{\zeta|\xi}=\deleted{\unsure{\beta}}\eta_0\lambda_1^\nu,
\qquad\text{i.e. \unsure{side} of (21)}
\]
\[
\tag{26bis}\unsure{\zeta|\xi}=
\]
where $\unsure{\zeta|\xi},\eta_0$ are two sections of $\underline M$ which
determine one another (for given $\nu$).
We have thus expressed in a
convenient way the loop equation, i.e. the
condition c)\inserted{of (17)}, it remains to express a) and b),
which are also formulated as
\[
\tag{27}
\left\{\begin{aligned}
\text{a)}\;&\unsure{\zeta|\xi}\ \text{locally of the form}\quad
\overset{\inserted{(\nu_3=0)}}{\alpha_0\rho(\alpha_0)^{-1}}
\quad\text{or}\quad
\overset{\inserted{(\nu_3=1)}}{\alpha_0\lambda_1^{-1}\rho(\alpha_0)^{-1}}\\
&\text{with }\alpha_0\in\Gamma\underline M\\
\text{b)}\;&\eta_0\ \text{locally of the form}\quad\beta_0\sigma(\beta_0)^{-1}
\end{aligned}\right.
\]
(Conditions whose formulation is independent
of $\nu,\nu_2$ \deleted{\unsure{for|$\nu_3$})}, but depends \unsure{apparently} on
$\nu_3$.) To make these conditions explicit, we
are going to write $\underline M$ additively, the relations
read
\margin{We have thus also set $\alpha=\alpha_0\tau^{\nu_3}$}
\[
\tag{28}
\left\{\begin{aligned}
\unsure{\zeta|\xi}&=\alpha_0-\rho(\alpha_0)
\quad\overset{\inserted{\text{case}}}{(\nu_3=0)}
\quad\text{or}\quad
\unsure{\zeta|\xi}=\alpha_0-\rho(\alpha_0)-\lambda_1
\quad\overset{\inserted{\text{case}}}{(\nu_3=1)}\\
\eta_0&=\beta_0-\sigma(\beta_0),
\end{aligned}\right.
\]
\margin{38}
\margin{38\\\texttt{https://grothendieck.umontpellier.fr}\\Université de Montpellier, UMR CNRS 5149, Case courrier 051, Place Eugène Bataillon\\34095 Montpellier cedex 5 - France}

\scanpage{40}

which we also express by
\[
\tag{29}
\left\{\begin{aligned}
\unsure{\zeta|\xi}&\in\operatorname{Im}(\mathrm{Id}_{\underline M}-\rho_{\underline M})
\quad(\nu_3=0)\quad\text{or}\quad
\unsure{\zeta|\xi}\in-\lambda_1+\operatorname{Im}(\mathrm{Id}_{\underline M}-\rho_{\underline M})\\
\eta_0&\in\operatorname{Im}(\mathrm{Id}_{\underline M}-\sigma_{\underline M}).
\end{aligned}\right.
\]

It is time to move on to explicit
matrix calculations, by referring $\underline M$ to
the basis $\lambda_0,\lambda_1$. In this basis, the
\unsure{operations} of $\rho_{\underline M},\sigma_{\underline M},\tau_{\underline M}$ are given by
\[
\tag{30}
\left\{\begin{aligned}
\rho_{\underline M}&=\begin{pmatrix}0&-1\\1&-1\end{pmatrix}
&\mathrm{id}_{\underline M}-\rho_{\underline M}&=\begin{pmatrix}1&1\\-1&2\end{pmatrix}\\
\sigma_{\underline M}&=\begin{pmatrix}0&1\\1&0\end{pmatrix}
&\mathrm{id}_{\underline M}-\sigma_{\underline M}&=\begin{pmatrix}1&-1\\-1&1\end{pmatrix}\\
\tau_{\underline M}&=\begin{pmatrix}-1&0\\0&-1\end{pmatrix}.&&
\end{aligned}\right.
\]

We have
\[
\det(\mathrm{id}_{\underline M}-\rho_{\underline M})=3,
\qquad\det(\mathrm{id}_{\underline M}-\sigma_{\underline M})=0
\]
which shows that \deleted{$\mathrm{id}_{M}$} the endomorphisms
$\mathrm{id}_{\underline M}-\rho_{\underline M}$ and $\mathrm{id}_{\underline M}-\sigma_{\underline M}$ of $\underline M$ are
not quite automorphisms\inserted{i.e. epimorphisms} ---
It screws up in characteristic 3 for the first, \uline{everywhere}
\deleted{\unsure{characteristic 2}}\inserted{(in every characteristic)} for the second.
\deleted{We shall therefore take}\inserted{We shall have}
\deleted{more} reason to take as ``parameters''
not $\unsure{\zeta|\xi}$ and $\eta_0$, but quite obviously
$\alpha_0$ and $\beta_0$ (in addition to the discrete ``parameters''
$\nu_2,\nu_3$ and $\nu'$), and to \deleted{define}\inserted{make} \deleted{thus} a
group (if possible\ldots) with the systems

\margin{39}
\margin{39\\\texttt{https://grothendieck.umontpellier.fr}\\Université de Montpellier, UMR CNRS 5149, Case courrier 051, Place Eugène Bataillon\\34095 Montpellier cedex 5 - France}

\scanpage{41}
\mspage{668}

\deleted{\unsure{(24)}}
\[
\tag{31}
\widetilde{\underline Z}_{\rho,\sigma}
=\left\{(\nu,\nu',\nu_2,\nu_3,\alpha_0,\beta_0)\ \left|\
\begin{gathered}
\nu,\nu'\in\Gamma\mathbf G_a,\quad
\nu_2,\nu_3\in\Gamma\{0,1\}_{S_0},\quad
\alpha_0,\beta_0\in\Gamma\underline M\\
\left\{\begin{aligned}
\unsure{\nu}&=2\nu'\unsure{-|+}\nu_2\qquad 2\nu+1\in\Gamma\mathbf G_m\\
\underbrace{\alpha_0-\rho(\alpha_0)\unsure{-|+}\nu_3\lambda_1}_{\unsure{\zeta|\xi}}
&=\underbrace{\beta_0-\sigma(\beta_0)}_{\eta}+\nu\lambda_1
\end{aligned}\right.
\end{gathered}\right.\right\}
\]
\margin{\uline{N.B.} $\nu'$ determines $\nu$ which is redundant}

The loop equation, in $\alpha_0,\beta_0$, thus takes
the form
\[
\tag{32}
(\unsure{\mathrm{id}|1}-\rho_{\underline M})\alpha_0
\unsure{-|=}(\mathrm{id}-\sigma_{\underline M})(\beta_0)
=\bigl(\underbrace{2\nu'\unsure{-|+}\nu_2}_{\nu}\unsure{+|-}\nu_3\bigr)\lambda_1
\]

\margin{This must indeed \unsure{come from the} fact that the homomorphism
\diagram{\begin{tikzcd}[column sep=large]
\underline M\times\underline M\arrow[r]&\underline M\\
(\alpha_0,\beta_0)\arrow[r,mapsto]&
(\mathrm{id}-\rho_{\underline M})\alpha_0-
(\deleted{\unsure{\mathrm{id}}}\inserted{1}-\sigma_{\underline M})(\beta_0)
\end{tikzcd}}
\unsure{is|is} indeed an epimorphism\ldots\ But this is the case only in characteristic $\ne3$ !}

\uline{Remark} The scheme $\underline Z_{\rho,\sigma}$ itself, as it
was defined in \S\ 49 \deleted{\unsure{\S\ 6}} $\overline{\underline{\mathrm{VI}}}$, is not smooth
it seems to me --- \unsure{it gets stuck|gets complicated} in characteristics 2 and 3.
On the other hand the finer group $\widetilde{\underline Z}_{\rho,\sigma}$,
defined by bringing in the
$\alpha,\beta$ right away (instead of $(\unsure{\zeta|\xi},\eta)$)
does seem \uline{smooth}. \unsure{This} leads once again
to modifying a little the presentation of the
constructions of \S\ 49, by \unsure{denying myself|not using}
\unsure{yet another} group $\underline Z_{\rho,\sigma}$ (tied to the coordinates
$\unsure{\zeta|\xi},\eta$), and by replacing it with
a group which \unsure{will appear|appears} \inserted{\unsure{at the same time|in a way}} as a
normalizer of $L_0$ (or of $L_0^{\unsure{\sharp}}$) in $\mathcal M_{\rho,\sigma}$
(or $\mathcal M_{\rho,\sigma}^{\unsure{\natural|\sharp}}$). We shall give a presentation along
these lines in the \deleted{\unsure{sequel}} \S\ that follows, and we
\unsure{assume} here that this construction is already in hand,



\margin{40}
\margin{40\\\texttt{https://grothendieck.umontpellier.fr}\\Université de Montpellier, UMR CNRS 5149, Case courrier 051, Place Eugène Bataillon\\34095 Montpellier cedex 5 - France}

\scanpage{42}

Whence the \deleted{\unsure{sphere|scheme}} (scheme \deleted{and) group}
$\widetilde{\underline Z}_{\rho,\sigma}$ will thus have
a suitable group structure, and
operates on $\unsure{\underline C^{+}|\underline G^{+}}$ via
\[
\tag{33}
\left\{\begin{aligned}
u(\rho)&=(\alpha_0-\rho(\alpha_0)+\nu_3\lambda_1)\rho\\
u(\sigma)&=(\beta_0-\sigma(\beta_0))\omega^{\nu_2}\sigma.
\end{aligned}\right.
\]

But we must all the same check that
these formulas (33) really do
determine an automorphism $u$ of $\unsure{\underline C^{+}|\underline G^{+}}$. The
essential check : make sure
that the \deleted{formula \unsure{[7]}}\inserted{relation} (7)
\[
\tag{34}[\sigma,\lambda_\infty]
\]
is respected by $u$.

I am going to make explicit the action of $u$ on
$\underline M$, by making explicit its values on
\[
\left\{\begin{aligned}
\lambda_0&=\omega l_0=\omega\sigma\rho\sigma\rho\\
\lambda_1&=\omega l_1=\omega\rho\sigma\rho\sigma
\end{aligned}\right.
\]
in terms of \deleted{$\unsure{\zeta|\xi}$ and $\eta_0$, \unsure{intro}}
\[
\tag{35}
\left\{\begin{aligned}
\unsure{\zeta|\xi}&=\alpha_0-\rho(\alpha_0)+\nu_3\lambda_1\\
\eta_0&=\beta_0-\sigma(\beta_0)
\end{aligned}\right.
\qquad\unsure{\alpha|\&}\qquad\in\Gamma\underline M
\]
Linked to $u$ by
\[
\tag{36}
u\rho=\unsure{\zeta|\xi}\rho,\qquad
u\sigma=\omega^{\nu_2}\eta_0\sigma.
\]
We shall have \deleted{thus} by the loop relation
\[
u(\varepsilon_0)=\deleted{\unsure{\mu}}\varepsilon_0^{\deleted{\unsure{\mu}}\inserted{(\mu)}}
\qquad(\mu=2\nu+1=4\nu'+2\nu_2+1)
\]
whence

\margin{41}
\margin{41\\\texttt{https://grothendieck.umontpellier.fr}\\Université de Montpellier, UMR CNRS 5149, Case courrier 051, Place Eugène Bataillon\\34095 Montpellier cedex 5 - France}

\scanpage{43}
\mspage{669}
\[
u(l_0)=u(\varepsilon_0^2)=l_0^{(\mu)}
\]
and
\[
u(\lambda_0)=\unsure{u}(\omega l_0)
=\underbrace{u(\omega)}_{=\omega}\,u(l_0)
=\omega l_0^{(\mu)}
=\underbrace{\omega l_0}_{\lambda_0}
 \bigl(\underbrace{l_0^{(2\nu)}}_{\lambda_0^{2\nu}}\bigr)
=\lambda_0^{(2\nu+1)}=\lambda_0^{(\mu)}
\]
as it should. (To reassure myself, I have checked that the second formula (33) does indeed give
\[
u(\sigma^2)=(\eta_0\sigma)(\eta_0\sigma)
=\bigl(\underbrace{\eta_0\,\sigma(\eta_0)}_{1}\bigr)
 \deleted{\unsure{\sigma^2}}\omega=\omega
\quad\text{i.e.}\quad\deleted{\text{and}}
\]
\[
\tag{37} u(\omega)=\omega
\]
We shall have
\[
u(\lambda_1)=u(\rho(\lambda_0))
=(\unsure{\zeta|\xi}\rho u)(\lambda_0)
=(\unsure{\zeta|\xi}\rho)u(\lambda_0)
=(\unsure{\zeta|\xi}\rho)(\lambda_0^\mu)=\lambda_1^\mu
\]
thus finally, as predicted
\[
\tag{38} u(\lambda)=\mu.\lambda\qquad\text{for}\quad\lambda\in\Gamma\underline M
\]
\margin{where $\mu=2\nu+1=4\nu'+2\nu_2+1$}
We shall indeed have that $u(\sigma)=(\beta_0-\sigma(\beta_0))\omega^{\nu_2}.\sigma$
commutes with $u(\lambda_\infty)=\mu.\lambda_\infty$.

All these calculations are a bit \unsure{murky}, logically, for at the same time I want to prove
the existence of a \deleted{\unsure{homomorphism}} $u$ \inserted{of $\unsure{\underline C^{+}|\underline G^{+}}$},
and in the preceding calculations I already assume it (\unsure{implicitly}, in a
``bigger'' group, namely the \unsure{smooth|free} group scheme \inserted{(\unsure{heuristic})}
generated by $\rho,\sigma\ldots$.
One way of making them \uline{rigorous} would be to use the intrinsic description of
$\unsure{\underline C^{+}|\underline G^{+}}$ :
\deleted{to \unsure{prove}}
\[
\tag{39}
\unsure{\underline C^{+}|\underline G^{+}}
=\unsure{C^{+}_{S_0}|G^{+}_{S_0}}
 \mathop{\wedge}_{M_{S_0}}\underline M_{S_0}
\]
thus defining a homomorphism\inserted{$u$} from
$\unsure{\underline C^{+}|\underline G^{+}}$ to an $S_0$-group $G$
(here $G=\unsure{\underline C^{+}|\underline G^{+}}$) amounts to
\margin{42}
\margin{42\\\texttt{https://grothendieck.umontpellier.fr}\\Université de Montpellier, UMR CNRS 5149, Case courrier 051, Place Eugène Bataillon\\34095 Montpellier cedex 5 - France}

\scanpage{44}
define homomorphisms
\[
\tag{40}
\diagram{\begin{tikzcd}[column sep=large]
\underline M_{S_0}\arrow[r,"u_{\underline M}"] & \unsure{G^{\sharp}|G^{\natural}}
\end{tikzcd}}
,\quad
\diagram{\begin{tikzcd}[column sep=large]
\unsure{C_{S_0}|G_{S_0}}\arrow[r,"u_{\unsure{C|G}}"] & \unsure{G^{\sharp}|G^{\natural}}
\end{tikzcd}}
\quad\text{i.e.}\quad
\diagram{\begin{tikzcd}[column sep=large]
\unsure{C|G}\arrow[r] & \unsure{G^{\sharp}(S_0)|G^{\natural}(S_0)}
\end{tikzcd}}
\]
satisfying the following two conditions:

a) Commutativity of the diagram
\[
\tag{41}
\diagram{\begin{tikzcd}[column sep=large,row sep=large]
M_{S_0}\arrow[r]\arrow[d] & \unsure{C^+_{S_0}|G^+_{S_0}}\arrow[ddr,"u_{\unsure{C|G}}"] & \\
\underline M_{S_0}\arrow[drr,"u_{\underline M}"'] & & \\
& & G
\end{tikzcd}}
\quad,
\]
and b)
\[
\tag{42}
\begin{gathered}
u_{\underline M}\text{ is compatible with the actions}\\
\text{of }\unsure{C^+_{S_0}|G^+_{S_0}}\text{ on }\underline M_{S_0}
\text{ (via its action on }M_{S_0}\text{)}\\
\text{and on }G\text{ (via }u_{\unsure{C|G}}\text{).}
\end{gathered}
\]
In the present case, we will have $G=\unsure{\underline C|\underline G}$
\[
\tag{43}
u_{\underline M}(\lambda)=\deleted{\unsure{\lambda}}\,i(\mu.\lambda)
\]
where
\[
i:\quad
\diagram{\begin{tikzcd}[column sep=large]
\underline M\arrow[r,hook] & \unsure{\underline C^+|\underline G^+}
\end{tikzcd}}
\]
is the inclusion, and
\[
\tag{44}
u_{\unsure{C|G}}:\deleted{\text{defined by}}
\inserted{\diagram{\begin{tikzcd}[column sep=large]
\unsure{C^+_{S_0}|G^+_{S_0}}\arrow[r]&G
\end{tikzcd}}}
\qquad\text{defined by (33),}
\]
We must therefore first check that (33), does define a homomorphism
\diagram{\begin{tikzcd}[column sep=small]\unsure{C_{S_0}|G_{S_0}}\arrow[r]&G\end{tikzcd}},
or what amounts to the same,
\diagram{\begin{tikzcd}[column sep=small]\unsure{C|G}\arrow[r]&G(S_0)\end{tikzcd}}.
\margin{43}
\margin{43\\\texttt{https://grothendieck.umontpellier.fr}\\Université de Montpellier, UMR CNRS 5149, Case courrier 051, Place Eugène Bataillon\\34095 Montpellier cedex 5 - France}

\scanpage{45}
\mspage{\textcircled{670}}
For this we \deleted{\unsure{note that} $u$ defines a homomorphism \unsure{to the} \unsure{of index} of $3$
$\unsure{C|G}$}
\inserted{use the presentation \deleted{\unsure{of}} of $\unsure{C^+|G^+}$}
\[
\deleted{\begin{gathered}
\unsure{C_0^+|G_0^+}\\
\unsure{\simeq}\\
\unsure{\operatorname{SL}(2,\mathbb Z)}=
\{\rho,\sigma\mid\rho^6=\sigma^4=1\}
\end{gathered}}
\]
\[
\tag{44}
\unsure{C^+|G^+}=
\left\{\rho,\sigma\ \middle|\ 
\begin{gathered}
\rho^6=\sigma^4=1,\quad\rho^3=\sigma^2\\
[\sigma,\lambda_\infty]=1\quad\text{where}\quad
\lambda_\infty=\deleted{\unsure{\rho}}\rho\sigma\rho^{-1}\sigma\rho
\end{gathered}
\right\},
\]
thus setting
\[
\tag{45}
\left\{\begin{aligned}
\rho'&=\unsure{\zeta\rho|u(\rho)},\quad
\sigma'=\eta_0\omega^{\nu_2}\sigma\qquad\text{cf. (33) and}\\
\unsure{\zeta|\xi}&=\alpha_0-\rho(\alpha_0)+\nu_3\lambda_1,\quad
\unsure{\eta|\eta_0}=\beta_0-\sigma(\beta_0)
\end{aligned}\right.
\]
we must check the relations
\[
\tag{46}
\rho'{}^6=\sigma'{}^4=[\sigma',\lambda'_\infty]=1,\quad
\rho'{}^3=\sigma'{}^2
\quad\text{where}\quad
\lambda'_\infty=\deleted{\unsure{\rho'(\lambda'_0)}}
\inserted{\rho'\sigma'\rho'{}^{-1}\sigma'\rho'}.
\]
We will begin by checking \deleted{\unsure{that|more}}
\[
\tag{47}\rho'{}^3=\sigma'{}^2=\omega
\]
which implies \deleted{\unsure{that}}
\[
\tag{48}
\rho'{}^6=\sigma'{}^4=1,\quad\text{thus}\quad
\exists\ \deleted{\text{unique}}
\diagram{\begin{tikzcd}[column sep=large]
\unsure{C_0^+|G_0^+}\arrow[r,"u_{\unsure{C_0|G_0}}"] & G(S_0)\\
\rho\arrow[r,mapsto]&\rho',\quad\sigma\mapsto\sigma'
\end{tikzcd}}
\]
Next, \deleted{setting} noting that we will have
\[
\lambda'_\infty=u_{\unsure{C_0|G_0}}(\lambda_\infty),
\]
it will suffice to \deleted{\unsure{this}} check
\[
\tag{49}\lambda'_\infty=\mu\lambda_\infty
\]
it will be obvious that $\lambda'_\infty$ commutes with
\[
\sigma'=u(\sigma)=\eta_0\sigma\qquad\deleted{\text{\unsure{pui}}}
\]
since $\sigma$ and $\eta_0\in M$ commute with it. Thus to construct
$u_{\unsure{C|G}}$, the essential checks are (47) and (49)
--- moreover once (47) is proved, the preceding calculations leading to (38)
are justified (\unsure{at least} where they make sense, i.e.,
for $u_{\unsure{C|G}}$ and $\lambda\in M$),
\margin{44}
\margin{44\\\texttt{https://grothendieck.umontpellier.fr}\\Université de Montpellier, UMR CNRS 5149, Case courrier 051, Place Eugène Bataillon\\34095 Montpellier cedex 5 - France}

\scanpage{46}
so (49) will indeed be verified.

We \deleted{will have} have
\[
\begin{aligned}
&\deleted{\unsure{u(\rho)=}}\\
\rho'{}^3&=\bigl(\unsure{\zeta|\xi}\rho(\unsure{\zeta|\xi})\rho^2(\unsure{\zeta|\xi})\bigr)\rho^3
=\unsure{\zeta|\xi}\rho(\unsure{\zeta|\xi})\rho^2(\unsure{\zeta|\xi})\,\omega\\
\sigma'{}^2&=(\eta_0\sigma(\eta_0))\sigma^2=(\eta_0\sigma(\eta_0))\omega
\end{aligned}
\]
so \deleted{and} formulas (47) reduce to:
\[
\tag{50}
\unsure{\zeta|\xi}\rho(\unsure{\zeta|\xi})\rho^2(\unsure{\zeta|\xi})=1,
\qquad\eta_0\sigma(\eta_0)=1
\]
i.e., additively
\[
\unsure{\zeta|\xi}+\rho(\unsure{\zeta|\xi})+\rho^2(\unsure{\zeta|\xi})=0,
\qquad\eta_0+\sigma(\eta_0)=0.
\]
whose verification is trivial, given that
\[
\rho_{\underline M}^3=\sigma_{\underline M}^2=\mathrm{id}_{\underline M},
\qquad(1+\rho_{\underline M}+\rho_{\underline M}^2)
=\deleted{\unsure{\mathrm{id}_{\underline M}}}0_{\underline M}.
\]
We have thus indeed constructed a homomorphism
\deleted{\unsure{to the}}
\[
u_{\unsure{C_0^+|G_0^+}}:\quad
\diagram{\begin{tikzcd}[column sep=large]
\unsure{C_0^+|G_0^+}\arrow[r]&\unsure{\underline C_0^+(S_0)|\underline G_0^+(S_0)}
\end{tikzcd}},
\]
and it remains to check (41) and (42).

For (41), we have already shown it at the \inserted{level of the existence of a $u_{\unsure{C_0|G_0}}$}
--- it amounts to proving
\[
\bigl(\quad u_{\unsure{C_0|G_0}}(\lambda_0)=\deleted{\unsure{\lambda_0}}\mu\lambda_0,
\quad u_{\unsure{C_0|G_0}}(\lambda_1)=\mu\lambda_1
\]
I will all the same redo the verification, to be sure.
\[
\begin{aligned}
\varepsilon'_0&=u_{\unsure{C_0|G_0}}(\underbrace{\varepsilon_0}_{\sigma\rho})
=u_{\unsure{C_0|G_0}}(\sigma)u_{\unsure{C_0|G_0}}(\rho)
=\eta_0\omega^{\nu_2}\deleted{\unsure{\zeta|\xi}}\sigma\unsure{\zeta|\xi}\rho
=\eta_0\omega^{\nu_2}\sigma(\unsure{\zeta|\xi})\sigma\rho\\
&=\eta_0\sigma(\unsure{\zeta|\xi})\omega^{\nu_2}\varepsilon_0
\end{aligned}
\]
\margin{45}
\margin{45\\\texttt{https://grothendieck.umontpellier.fr}\\Université de Montpellier, UMR CNRS 5149, Case courrier 051, Place Eugène Bataillon\\34095 Montpellier cedex 5 - France}

\scanpage{47}
\mspage{\unsure{671}}

I claim that
\[
\begin{aligned}
\varepsilon'_0&=\varepsilon_0^{(\mu)}\quad\text{i.e.}\quad
\eta_0\sigma(\unsure{\zeta|\xi})\omega^{(\nu_2)}
 =\underbrace{\varepsilon_0^{(\mu-1)}}_{l_0^\nu}
\quad\text{where }\mu-1=2\nu\deleted{\unsure{+1|=1}},\quad\text{i.e.}\\
\eta_0\sigma(\unsure{\zeta|\xi})\omega^{\nu_2}&=l_0^{(\nu)}
\quad\text{where }\nu=2\nu'\unsure{-|+}\nu_2\quad\text{i.e.}\\
\eta_0\sigma(\unsure{\zeta|\xi})
 &=\underbrace{l_0^{(2\nu')}}_{\lambda_0^{2\nu'}}
 \underbrace{\unsure{(\omega l_0)^{(-\nu_2)}|(\omega l_0^{(-\nu_2)})}}_{\lambda_0^{-\nu_2}}
 =\underbrace{\lambda_0^{(2\nu'\unsure{-|+}\nu_2)}}_{\lambda_0^\nu}
 \quad\text{i.e.}\\
\eta_0\sigma(\unsure{\zeta|\xi})&=\lambda_0^{(\nu)}
\end{aligned}
\]
i.e., applying $\sigma$ and noting that
$\sigma(\eta_0)=\eta_0^{-1}$, $\sigma^2(\unsure{\zeta|\xi})=\unsure{\zeta|\xi}$,
$\sigma(\lambda_0)=\lambda_1$,
\[
\eta_0^{-1}\unsure{\zeta|\xi}=\lambda_1(\nu)
\quad\text{i.e. finally }\deleted{\unsure{\zeta_0=|\xi_0=}}
\unsure{\zeta|\xi}=\eta_0\unsure{\lambda_1(\nu)|\lambda_1^{(\nu)}}
\quad\text{or}
\]
\[
\unsure{\zeta|\xi}=\eta_0+\nu\lambda_1\quad\text{(written additively)}
\]
which is true by hypothesis (cf. (31)).

We will then have
(\deleted{\unsure{for $u=u_{\unsure{C_0|G_0}}$|for $u_{\unsure{C_0|G_0}}$}}
\inserted{writing $u$ instead of $u_{\unsure{C_0|G_0}}$ to simplify})
\[
u_{\deleted{\unsure{C_0|G_0}}}(\lambda_1)
 =u\rho(\lambda_0)
 =(\unsure{\zeta|\xi}\rho u)(\lambda_0)
 =\unsure{\zeta|\xi}\rho(u(\lambda_0))
 =\unsure{\zeta|\xi}\rho(\lambda_0(\mu))
 =\lambda_1(\mu)
\]
OK.

It remains to check (42), which \deleted{\unsure{means}}
\inserted{amounts} to the \unsure{commutativity} of
$u_{\underline M}=\mu\,\mathrm{id}_{\underline M}$ with
$(\rho_{\underline M},\rho'_{\underline M})$,
$(\sigma_{\underline M},\sigma'_{\underline M})$
--- \unsure{this means} also that
$\rho_{\underline M}$ and $\rho'_{\underline M}$ on the one hand,
$\sigma_{\underline M}$ and $\sigma'_{\underline M}$ on the other hand,
\deleted{\illegible{2}}\inserted{are equal, which}
is also obvious.

Thus, the group $\widetilde Z_{\rho,\sigma}$ is well and
truly constructed --- at least, as a
scheme over $S_0=\mathrm{Spec}\,\mathbb Z$, with morphisms
of schemes
\[
\tag{51}
\diagram{$
\begin{tikzcd}[column sep=3em,row sep=.25em]
\widetilde{\underline Z}_{\rho,\sigma}\arrow[r,"\nu'"] & G_a\\
{}\arrow[r,"\deleted{\unsure{\nu|\varepsilon}}"] & \deleted{\unsure{\{\pm1\}_{S_0}^{2}|\mu_{S_0}^{2}}}\\
{}\arrow[r,"\nu_2"] & \{0,1\}_{S_0}\\
{}\arrow[r,"\nu_3"] & \{0,1\}_{S_0}
\end{tikzcd}
\quad\Bigg|\quad
\begin{gathered}
\text{whence }\mu=4\nu'+2+\nu_2:\quad
\begin{tikzcd}[column sep=1.5em,row sep=1em]
\widetilde Z_{\rho,\sigma}\arrow[r]\arrow[rd] & G_a\\
 & G_m\arrow[u,hook]
\end{tikzcd}\\
\varepsilon_2,\varepsilon_3:\quad
\begin{tikzcd}[column sep=3em]
\widetilde Z_{\rho,\sigma}\arrow[r,shift left=.3ex]\arrow[r,shift right=.3ex]
 & \mu_{S_0}=\{\pm1\}_{S_0}
\end{tikzcd}
\end{gathered}
$}
\]

\margin{46}
\margin{46\\\texttt{https://grothendieck.umontpellier.fr}\\Université de Montpellier, UMR CNRS 5149, Case courrier 051, Place Eugène Bataillon\\34095 Montpellier cedex 5 - France}

\scanpage{48}

and a morphism of \deleted{\unsure{functors}} \inserted{presheaves} over $S_0$
\[
(52)\qquad
\diagram{\begin{tikzcd}[row sep=small]
\widetilde{\underline Z}_{\rho,\sigma} \arrow[r] & \underline{\operatorname{Aut}}_{\mathrm{gr}}\bigl(\unsure{\underline C^{+}|\underline G^{+}}\bigr),\\
\underline u \arrow[r,mapsto] & u
\end{tikzcd}}
\]
giving rise to a morphism
\[
(53)\qquad
\diagram{\begin{tikzcd}[row sep=small]
\widetilde{\underline Z}_{\rho,\sigma} \arrow[r] & \underline{\operatorname{Aut}}_{\mathrm{gr}}\bigl(\unsure{\underline C^{+}|\underline G^{+}}\bigr)\times\mathbf G_{\mathrm m}\times\mu_{S_0}\times\mu_{S_0}\\
\underline u \arrow[r,mapsto] & u,\mu(\underline u),\varepsilon_2(\underline u),\varepsilon_3(\underline u)
\end{tikzcd}}
\]
--- calculations already done thirty-six times in various contexts should show that this morphism (which is \uline{monomorphic}\ \unsure{?}\ in characteristic $\ne 2,3$) is a morphism of \uline{groups}, for a suitable \uline{\unsure{group|groups}} structure on $\widetilde{\underline Z}_{\rho,\sigma}$, uniquely determined (since it is so in characteristic $\ne 2,3$, and we conclude by reason of schematic density). We will have \deleted{\unsure{this is not}}
\[
(54)\qquad
\diagram{\begin{tikzcd}
1\arrow[r]&\widetilde{\underline Z}_{\rho,\sigma}^{+0}\arrow[r]&\widetilde{\underline Z}_{\rho,\sigma}\arrow[r]&\mathbf G_{\mathrm m}\times\mu_{S_0}\times\mu_{S_0}\arrow[r]&1
\end{tikzcd}}
\]
\margin{a bit \unsure{cavalier}, since
\diagram{\begin{tikzcd}
\underline Z_{\rho,\sigma}\arrow[r]&\mathbf G_{\mathrm m}\times\mu_{S_0}\times\mu_{S_0}
\end{tikzcd}}
is not an epimorphism, because
\diagram{\begin{tikzcd}[row sep=small]
\underline M\times\underline M\arrow[r]&\underline M\\
(\alpha_0,\beta_0)\arrow[r,mapsto]&(1-\rho_{\underline M})\alpha_0\unsure{-|+}(1-\sigma_{\underline M})\beta_0
\end{tikzcd}}
is not epimorphic --- \unsure{where there} is a \unsure{swindle} in characteristic 3 --- cf. below\ldots}
making explicit
\[
(55)\qquad
\begin{gathered}
\widetilde{\underline Z}_{\rho,\sigma}^{+0}\underset{\inserted{\text{subgroup scheme}}}{\subset}\underline M\times\underline M\times\mathbf G_{\mathrm a}\\
\Vert\\
\left\{(\alpha_0,\beta_0,\nu')\,\middle|\,
\begin{gathered}
\deleted{\unsure{\alpha_0}}4\nu'=0\\
(1+\rho_{\underline M})\alpha_0=(1-\sigma_{\underline M})\beta_0+2\nu'\lambda_1
\end{gathered}\right\}
\end{gathered}
\]
We will thus have
\[
(56)\qquad
\diagram{\begin{tikzcd}
1\arrow[r]&\widetilde{\underline Z}_{\rho,\sigma}^{+00}\arrow[r]&\widetilde{\underline Z}_{\rho,\sigma}^{+0}\arrow[r]&\underline{{}_4\mathbf G_{\mathrm a}}\arrow[r]&1
\end{tikzcd}}
\]
\margin{47}
\margin{47\\\texttt{https://grothendieck.umontpellier.fr}\\Université de Montpellier, UMR CNRS 5149, Case courrier 051, Place Eugène Bataillon\\34095 Montpellier cedex 5 - France}

\scanpage{49}
\mspage{672}

where
\[
(57)\qquad
\begin{gathered}
\widetilde{\underline Z}_{\rho,\sigma}^{+00}\overset{\mathrm{dfn}}{=}
\left\{\underline u\in\Gamma\widetilde{\underline Z}_{\rho,\sigma}\,\middle|\,
\nu'(\underline u)=\nu_2(\underline u)=\nu_3(\unsure{u|\underline u})=0\right\}\\
\simeq\\
\left\{(\alpha_0,\beta_0)\in\Gamma(\underline M\times\underline M)\,\middle|\,
(1-\rho_{\underline M})\alpha_0=(1-\sigma_{\underline M})\beta_0\right\}.
\end{gathered}
\]
\margin{in characteristic 3 its relative dimension \unsure{jumps} from 2 to 3 !}
\uline{group} scheme $\underline Z_{\rho,\sigma}^{+00}$ is smooth \inserted{in characteristic $\ne 3$} of relative dimension 2 --- it is even canonically a Module scheme, we will have
\[
(58)\qquad\widetilde{\underline Z}_{\rho,\sigma}^{+00}\simeq\unsure{\mathbf V|\mathbb V}(\underline P)
\]
where $P$ is the free $\mathbb Z$-module of rank 2 defined \deleted{\illegible{1}} by the exact sequence
\[
(59)\qquad
\diagram{\begin{tikzcd}[row sep=small]
0\arrow[r]&\deleted{\unsure{M}}\inserted{P}\arrow[r]&M\times M\arrow[r]&M\arrow[r]&0\\
&&(\alpha_0,\beta_0)\arrow[r,mapsto]&(1-\rho_{\underline M})\alpha_0-(1-\sigma_{\underline M})\beta_0&
\end{tikzcd}}
\]
The homomorphism $M\times M\longrightarrow M$ is made explicit, setting
\[
(60)\qquad
\left\{\begin{aligned}
\alpha_0&=a_0\lambda_0+\unsure{a_1}\lambda_1
&\text{hence}\quad\alpha_0-\rho\alpha_0&=(a_0+\unsure{a_1})\lambda_0+(2\unsure{a_1|b_1}-a_0)\lambda_1\\
\beta_0&=b_0\lambda_0+b_1\lambda_1
&\text{hence}\quad\beta_0-\sigma\beta_0&=(b_0-b_1)\lambda_0+(b_1-b_0)\lambda_1
\end{aligned}\right.
\]
as
\[
(61)\qquad
\diagram{\begin{tikzcd}
(\alpha_0,\beta_0)\arrow[r,mapsto]&
\underbrace{(a_0+a_1+(b_1\unsure{-|+}b_0))}_{c_0}\lambda_0+
\underbrace{(2a_1-a_0-(b_1-b_0))}_{c_1}\lambda_1
\end{tikzcd}}
\]
\deleted{\unsure{We}} This is a surjective homomorphism \inserted{except in characteristic 3}, \deleted{for given $(c_0,c_1)$,}
\deleted{\unsure{we will set}\ $\unsure{a_1|c_1}$\ \illegible{1}}
\inserted{--- the elements of the image are the}
\[
(62)\qquad
\left\{c_0\lambda_0+c_1\lambda_1\,\middle|\,
\underbrace{c_0+c_1}_{=3a_1}\equiv0\quad(3)\right\}
\]
\uline{Remark} If we take \deleted{\unsure{a}}\inserted{a} $u$ defined over $\widehat{\mathbb Z}$ coming from an element of $M[0]$, then
\margin{48}
\margin{48\\\texttt{https://grothendieck.umontpellier.fr}\\Université de Montpellier, UMR CNRS 5149, Case courrier 051, Place Eugène Bataillon\\34095 Montpellier cedex 5 - France}

\scanpage{50}

we will have \deleted{that}
\[
\mu\in\widehat{\mathbb Z}^{*},\quad
\mu\equiv\deleted{\unsure{\varepsilon_3|\varepsilon_2}}(-1)^{\nu_3}(=1-2\nu_3)\quad\mathrm{mod}\ (3)
\]
and the condition on $(\mu,\varepsilon_2,\varepsilon_3)$ for being in the image of
\[
\diagram{\begin{tikzcd}
\widetilde{\underline Z}_{\rho,\sigma}\arrow[r]&\mathbf G_{\mathrm m}\times\mu_{S_0}\times\mu_{S_0}
\end{tikzcd}},
\]
which is written in general (by virtue of (62) \unsure{\&} and of (32))
\[
2\nu'-\nu_2+\nu_3\equiv0\quad(3)
\]
becomes, on multiplying by 2
\[
\underbrace{4\nu'-2\nu_2}_{\mu-1}+2\nu_3\equiv0\quad(3)
\]
i.e.
\[
\mu\equiv1-2\nu_3\quad(3),
\]
and it is automatically satisfied.

\deleted{
That is why I feel like modifying the definition of \deleted{$\unsure{Z}$} $\widetilde{\underline Z}_{\rho,\sigma}$, by defining $\underline Z_{\rho,\sigma}$ as
\[
\deleted{(63)\quad\widetilde{\underline Z}_{\rho,\sigma}=}
\left\{\begin{gathered}
\deleted{\unsure{\Gamma|\mu}}(\inserted{\mu,}\nu',\nu'',\nu_2,\nu_3,\alpha_0,\beta_0)
\in\mathbf G_{\mathrm m}\times\mathbf G_{\mathrm a}\times\unsure{\mathbf G_{\mathrm a}}\times\{0,1\}_{S_0}\times\{0,1\}_{S_0}\times\underline M\times\underline M\\
\left|\begin{gathered}
\mu=4\nu'\unsure{-|+}2\unsure{\nu_2}+1(=6\nu''-2\nu_3+1),\quad 2\nu'+\nu_3-\nu_2=3\nu'',\\
(\unsure{1}-\rho_{\underline M})(\alpha_0)-(1-\sigma_{\underline M})(\beta_0)
=\underbrace{(\deleted{\unsure{2\nu'-\nu_2+\nu_3}}\inserted{3\nu''})}_{c}\lambda_1
\end{gathered}\right.
\end{gathered}\right\}
\]
}
\margin{by \unsure{expressing|requiring} $\unsure{\mu|\nu}$\ \unsure{are written}
\[
\begin{aligned}
\mu&\equiv1-2\nu_2\quad(4)\\
\mu&\equiv1-2\nu_3\quad(6)
\end{aligned}
\unsure{\left[\vphantom{\begin{matrix}\mu\\\mu\end{matrix}}\right.}
\]
\deleted{the loop equation}
\[
\begin{aligned}
\mu&=1\unsure{-|+}2\nu_2+4\nu'&&\text{\unsure{with}}\\
\mu&=1-2\nu_3+6\nu''\\
-\nu_2+2\nu'&=-\nu_3+3\nu''&&\text{i.e.}\\
2\nu'+\nu_3-\nu_2&=3\nu''
\end{aligned}
\]}
\deleted{
\deleted{We will have here
\[
2c=6\nu'',\quad\text{and we \unsure{will|have} in fact}
\]
take $c=3\nu''$, and.}
We will have
\[
(63)\qquad
\begin{gathered}
\widetilde{\underline Z}_{\rho,\sigma}\subset\mathbf G_{\mathrm a}\times\mathbf G_{\mathrm a}\times\{0,1\}_{S_0}\times\{0,1\}_{S_0}\times\underline M\times\underline M\\
\Vert\\
\left\{(\nu',\nu'',\nu_2,\nu_3,\alpha_0,\beta_0)\,\middle|\,
\begin{aligned}
\text{a)}\quad &2\nu'\deleted{+\unsure{\nu_3}}-\nu_2=3\nu''-\nu_3\\
&\deleted{(\text{hence }4\nu'\unsure{-|+}2\nu_2+1=6\nu''-2\nu_3+1,\ \text{i.e. }\mu)}\\
&\inserted{\text{i.e. }\nu,\ \text{hence }\mu=2\nu+1}\\
\text{b)}\quad &\mu\in\Gamma\mathbf G_{\mathrm m}\\
\text{c)}\quad &\text{\unsure{loop|loop} equation}\\
&(1-\rho_{\underline M})(\alpha_0)-(1-\sigma_{\underline M})(\beta_0)=3\nu''\unsure{\underline{\lambda_1}|\lambda_1}
\end{aligned}\right\}
\end{gathered}
\]
}
\margin{49}
\margin{49\\\texttt{https://grothendieck.umontpellier.fr}\\Université de Montpellier, UMR CNRS 5149, Case courrier 051, Place Eugène Bataillon\\34095 Montpellier cedex 5 - France}

\scanpage{51}
\mspage{673}

3. The groups \deleted{$\unsure{\widetilde Z_{\rho,\sigma}}$ \unsure{revised} \illegible{2} $\mathbb O_N$} $\mathbb O_N^{*}$.

We are going to introduce an auxiliary construction,
by defining a group scheme which
corresponds to the datum of quantities
$\mu\in\Gamma\mathbf G_{\mathrm m}$, $\nu',\nu''\in\Gamma\mathbf G_{\mathrm a}$,
$\nu_2,\nu_3\in\Gamma\{0,1\}_{S_0}$, satisfying
\[
(63)\qquad
\mu=4\nu'-2\nu_2+1=6\nu''-2\nu_3+1,
\qquad 2\nu'-\nu_2=3\nu''-\nu_3.
\]

In general, let $N$\inserted{$\in\mathbb N^{*}$} be a fixed \inserted{natural} integer
(the case which interests us for what follows is
$N=12$), and let us consider the scheme \inserted{in groups, $\mathbb O_N$}
over $\mathbb Z$ associated to the extension of discrete groups
\deleted{\[(64)\qquad \mathbb O_N\]}
\[
(64)\qquad
\diagram{\begin{tikzcd}[column sep=large]
1\arrow[r]&\mathbb Z\arrow[r,"N"]&\mathbb Z\arrow[r]&\mathbb Z/N\mathbb Z\arrow[r]&1
\end{tikzcd}}
\]
which fits into an exact sequence
\[
(65)\qquad
\diagram{\begin{tikzcd}[column sep=large]
1\arrow[r]&\unsure{\mathbf E^{1}|\mathbb E^{1}}\arrow[r,"i"]&\mathbb O_N\arrow[r]&(\mathbb Z/N\mathbb Z)_{S_0}\arrow[r]&1.
\end{tikzcd}}
\]

Let us recall the \deleted{\unsure{defi}} simple construction
\[
(66)\qquad
\mathbb O_N
=(\mathbb Z)_{S_0}\mathop{\wedge}_{(N\mathbb Z)}\underline{N\mathbb Z}.
\]
\deleted{But} I say that there is \unsure{in|on} $\mathbb O_N$ a
unique Ring structure, whose
addition is the \inserted{abelian} group structure already
envisaged, and where multiplication is
obtained as follows: A section of $\mathbb O_N$ is
\deleted{local} of the form \deleted{$\unsure{k|x}$} $n+N[\lambda]$, with
the rules of calculation
\[
(67)\qquad
\left\{\begin{aligned}
\text{a)}\quad &(n+N[\lambda])+(n'+N[\lambda'])=(n+n')+N[\lambda+\lambda']\\
\text{b)}\quad &n+N[\lambda]=n'+N[\lambda']\Longleftrightarrow
\deleted{\unsure{\xi|n}}\ n+N\lambda=n'+N\lambda'\text{ in }\Gamma\mathbf G_{\mathrm a}.
\end{aligned}\right.
\]

\margin{50}
\margin{50\\\texttt{https://grothendieck.umontpellier.fr}\\Université de Montpellier, UMR CNRS 5149, Case courrier 051, Place Eugène Bataillon\\34095 Montpellier cedex 5 - France}

\scanpage{52}

We then multiply the sections by
\[
(68)\qquad
(\unsure{n}+N[\lambda])(n'+N[\lambda'])
=nn'+N[n\lambda'+n'\lambda+\lambda\lambda'],
\]
we check that it is \deleted{\unsure{biad}} compatible with
the relations (67 b), that it is biadditive,
commutative, associative, that $1+N[0]$ is
the unit\ldots{} Another way of \deleted{\unsure{it|it} \unsure{con|en}} seeing the
ring structure. \deleted{We} Let more
generally \deleted{$\mathbb O_N=\mathbb Z_{S_0}$}
$A$ be a \deleted{$\unsure{\mathbb Z|z}$} $\mathbb Z$-algebra,
\deleted{\unsure{com}} associative \inserted{\unsure{unital|unitary}} (say)
\deleted{\unsure{principal}} \inserted{equipped with an \uline{ideal}}
\margin{\inserted{two-sided}} $J$ \inserted{which is}
a free $\mathbb Z$-module of finite type, let us consider
\deleted{\unsure{an} ideal}
\deleted{the ideal $NA$} and the extension
structure
\[
(69)\qquad
\diagram{\begin{tikzcd}[column sep=large]
0\arrow[r]&\deleted{\unsure{NA}}\inserted{J}\arrow[r]&A\arrow[r]&\overbrace{A/\deleted{\unsure{NA|N}}\inserted{J}}^{B}\arrow[r]&0
\end{tikzcd}}
\]
hence a commutative group scheme
$\underline A_{[N]}$, which is obtained as
a quotient of
\deleted{$(A)_{\unsure{S|S_0}}\cdot\underline{\unsure{NB|NA}}\;\unsure{\xi\to|p}$}
semidirect \inserted{product $A_{S_0}\cdot\underline J$})
\[
(70)\qquad
\diagram{\begin{tikzcd}[column sep=large]
0\arrow[r]&\underline J\arrow[r]&A_{S_0}\times\underline J\arrow[r]&\underline A_N\arrow[r]&0\\
&\lambda\arrow[r,mapsto]&(\lambda,-\lambda)&&
\end{tikzcd}}
\]
and we find that $A_{S_0}\cdot\underline J$ has a
$\frac12$-direct product \uline{Ring} structure, and that

\margin{51}
\margin{51\\\texttt{https://grothendieck.umontpellier.fr}\\Université de Montpellier, UMR CNRS 5149, Case courrier 051, Place Eugène Bataillon\\34095 Montpellier cedex 5 - France}

\scanpage{53}
\mspage{674}

the image of $\underline J$ in it is an ideal ...

We do find a ring structure on $\underline A_{[N]}$, and an exact sequence \deleted{\unsure{of}}
\[
\tag{71}
\diagram{\begin{tikzcd}[column sep=large]
1 \arrow[r,no head] & \underline J \arrow[r,no head] & \underline A_{[N]} \arrow[r] & B_{S_0} \arrow[r] & 1
\end{tikzcd}}
\]
where $\underline A_{[N]}\mathbin{-}B_{S_0}$ is a \unsure{morphism} of \uline{ring schemes}. Passing to the multiplicative groups \deleted{\unsure{rings}}, we find an exact sequence
\[
\tag{72}
\diagram{\begin{tikzcd}[column sep=large,row sep=small]
1 \arrow[r] & (1+\underline J)^* \arrow[r,no head] & A_{[N]}^* \arrow[r,no head] & (B_{S_0})^* \arrow[r,no head] \arrow[d,equal] & 1 \\
& & & (B^*)_{S_0} &
\end{tikzcd}}
\]
where for $A_{[N]}^*$, $(B_{S_0})^{\unsure{*|'}}$ and $B^*$, the $*$ denotes the \unsure{passage} to the multiplicative \unsure{subgroup} and where $(1+\underline J)^*$ has the analogous meaning
\[
\tag{73}
\diagram{\begin{tikzcd}[row sep=large]
(1+\underline J)^*\unsure{\subset|\subseteq}1+\underline J \arrow[d,draw=none,"="{description,rotate=90}] \\
\{(1+\lambda)\mid 1+\lambda\ \text{invertible in }\underline A\} \arrow[d,draw=none,"\simeq"{description,rotate=90}] \\
\{(\lambda,\lambda')\in\underline J\times\underline J\mid\lambda+\lambda'+\lambda\lambda'=0\ \text{in }\underline J\}.
\end{tikzcd}}
\]

Thus we will have, in the special case $A=\mathbb Z$, $J=N\mathbb Z$
\[
\tag{74}
\diagram{\begin{tikzcd}[column sep=large]
1 \arrow[r] & (\underline{1+N\mathbb Z})^* \arrow[r] & \mathbb O_N^* \arrow[r,no head] & (\mathbb Z/N\mathbb Z)^*_{S_0} \arrow[r,no head] & 1
\end{tikzcd}}
\]
where the kernel

\margin{52}
\margin{52\\\texttt{https://grothendieck.umontpellier.fr}\\Université de Montpellier, UMR CNRS 5149, Case courrier 051, Place Eugène Bataillon\\34095 Montpellier cedex 5 - France}

\scanpage{54}
\deleted{\mspage{\unsure{674}}}

$(1+N\mathbb Z)^*$ is defined as
\[
\tag{75}
(1+N\mathbb Z)^*\overset{\mathrm{dfn}}{=}
\{\lambda\in\deleted{\unsure{\mathbf G_a}}\underline{\mathbb Z}=\mathbf G_a\mid 1+N\lambda\ \text{invertible}\}
\]

\deleted{We have a canonical \unsure{homomorphism} (\unsure{in part}}

The $\mathbb O_N$ for varying $N$ form a projective system of Rings, the $\mathbb O_N^*$ form a projective system of commutative group schemes.

I am not going to look more closely at the structure of these projective systems, nor examine \unsure{the} projective limits --- \unsure{questions} which arise more \unsure{general|generally} when one starts from an $A$ as above. We are going to make explicit the canonical \unsure{homomorphism}
\[
\tag{76}
\diagram{\begin{tikzcd}[column sep=large]
\mathbb O_N^* \arrow[r,no head] & \mathbb O_1^*=\mathbb O^*=\mathbf G_m
\end{tikzcd}}
\]
\deleted{$\unsure{\varphi}$}

\deleted{whose kernel is}
which in general is neither an epimorphism, nor a monomorphism. Its kernel is \deleted{formed of the} \deleted{$\unsure{\varphi}\ \unsure{\lambda}$} isomorphic to the \unsure{subgroup|system} of
\unsure{$+$}
\[
\tag{77}
\{1+N[\lambda]\mid 1+N\lambda=1\ \text{i.e.}\ N\lambda=0\}
\]
which as a scheme (\unsure{without} group structure) is isomorphic to ${}_N\mathbf G_a$,

\margin{53}
\margin{53\\\texttt{https://grothendieck.umontpellier.fr}\\Université de Montpellier, UMR CNRS 5149, Case courrier 051, Place Eugène Bataillon\\34095 Montpellier cedex 5 - France}

\scanpage{55}
\mspage{675}

which is a group scheme (not flat if $N\ne1$) which is zero over $\operatorname{Spec}\mathbb Z[1/N]$, and \deleted{is \unsure{dimension|of dimension} relative $1$} the fibers at the $p\mid N$ are $\simeq\mathbf G_a$.

I realize now that it is a clumsy move to leave \unsure{the} \unsure{group} $\mathbb O_N^*$, which is \unsure{perfect} of its kind, to refer to $\mathbf G_m$.

I am going to examine in particular
\[
\tag{78}
\diagram{\begin{tikzcd}[column sep=large]
1 \arrow[r,no head] & (1+12\mathbb Z)^* \arrow[r,no head] & \mathbb O_{12}^* \arrow[r] & (\mathbb Z/12\mathbb Z)^*_{S_0} \arrow[r,no head] & 1.
\end{tikzcd}}
\]

Let us note that \deleted{$\mathbb Z/12\mathbb Z$}
\[
(\mathbb Z/12\mathbb Z)^*\simeq(\mathbb Z/4\mathbb Z)^*\times(\mathbb Z/6\mathbb Z)^*\simeq\mu\times\mu.
\]

Thus (78) takes the form
\deleted{and we find a set-theoretic section \unsure{the|the}\ \illegible{1}\ \unsure{of|of the}\ \unsure{group|kernel}\ \unsure{over}\ of the extension $\unsure{\beta}$}
\[
\tag{79}
\diagram{\begin{tikzcd}[column sep=large]
1 \arrow[r,no head] & (1+12\mathbb Z)^* \arrow[r,no head] & \mathbb O_{12}^* \arrow[r,no head,"\varepsilon_2\times\varepsilon_3"] & \mu\times\mu \arrow[r,no head] & 1
\end{tikzcd}}
\]

The quantities $\nu_{\unsure{2|3}}\ \nu_3$ are also \inserted{morphisms} defined on $\mathbb O_{12}^*$
\margin{\deleted{\unsure{here}}}
\[
\tag{80}
\diagram{\begin{tikzcd}[column sep=large]
\mathbb O_{12}^* \arrow[r,"(\nu_2,\nu_3)"] & \{0,1\}_{S_0}\times\{0,1\}_{S_0}
\end{tikzcd}}
\]
they are the composites
\[
\tag{81}
\diagram{$\begin{tikzcd}[ampersand replacement=\&,column sep=large,row sep=small]
\mathbb O_{12}^* \arrow[r,"(\varepsilon_2,\varepsilon_3)"] \& \mu\times\mu \arrow[r,"\sim"] \& \{0,1\}_{S_0}\times\{0,1\}_{S_0} \\
\& \& \text{deduced from}\quad\deleted{\unsure{(+1)\mapsto(0)}}\quad\begin{array}{rcl}1&\mapsto&0\\-1&\mapsto&1\end{array}
\end{tikzcd}$}
\]

\margin{54}
\margin{54\\\texttt{https://grothendieck.umontpellier.fr}\\Université de Montpellier, UMR CNRS 5149, Case courrier 051, Place Eugène Bataillon\\34095 Montpellier cedex 5 - France}

\scanpage{56}

\deleted{\unsure{that}}

and they are group homomorphisms (for addition modulo $2$ on $\{0,1\}_{S_0}$.

\inserted{We also have $\varepsilon_2=(-\unsure{1})^{\nu_2}=1-2\nu_2,\quad\varepsilon_3=(-\unsure{1})^{\nu_3}=1-2\nu_3$}

The ``four cases of \uline{signs}'' in $(\mathbb Z/12\mathbb Z)^*$ are given in the table
\[
\tag{82}
\diagram{$\begin{array}{r|c|c}
{}_{\varepsilon_3}\!\diagdown^{\varepsilon_2}
&\begin{array}{c}+1\\\text{i.e. }\nu_2=0\end{array}
&\begin{array}{c}-1\\\text{i.e. }\nu_2=1\end{array}\\ \hline
\begin{array}{r}+1\\\text{i.e. }\nu_3=0\end{array}&1&+5\\ \hline
\begin{array}{r}-1\\\text{i.e. }\nu_3=1\end{array}&-5&-1\\ \hline
\end{array}$}
\]

thus we find a set-theoretic section of (79) by taking the sections of $\underline{\mathbb O}_{12}^*$ of the form
\[
\tag{83}
\underline{\mu}_0=(-5)^{\nu_2}(+5)^{\nu_3}
\left\{\begin{array}{ll}
\equiv(-1)^{\nu_3}&(6)\\
\equiv(-1)^{\nu_2}&(4)
\end{array}\right.
\]

Thus a general section of $\mathbb O_{12}^*$ is written
\[
\tag{84}
\underline{\mu}=\underbrace{(-5)^{\nu_2}(+5)^{\nu_3}}_{\mu_0}+\unsure{12}[\unsure{\deleted{\lambda_0}\inserted{\nu_0}|\nu_0}]
\]
where
\[
\left\{\begin{array}{l}
(\nu_2,\nu_3)\in\Gamma(\{0,1\}_{S_0}\times\{0,1\}_{S_0})\\
\lambda_0\in\Gamma(\mathbf G_a).
\end{array}\right.
\]

\deleted{We \unsure{argue}\ \unsure{by|more}\ \unsure{modulo}\ $4$}
\inserted{Let us note that we have}
\[
\deleted{\mu=(4+1}
\qquad
5^{\nu_3}=1+4\nu_3
\]
and
\[
(-5)^{\nu_2}=\left[\begin{aligned}
(-1)^{\nu_2}5^{\nu_2}&=(1-2\nu_2)(1+4\nu_2)\\
&=1+2\nu_2-8\nu_2^2
\end{aligned}\right]=1-6\nu_2
\]
\margin{\uline{N.B.} $\nu_2=\nu_2^2$}

thus
\[
\tag{\unsure{85}}
\begin{aligned}
\underline{\mu}_0&=(1+4\nu_3)(1\unsure{-6}\nu_2)\deleted{\unsure{-8}\nu_2^2)}\\
&\deleted{=1+2\nu_2+\unsure{(-8)}\nu_2^2+4\nu_3(1+2\nu_2-8\nu_2^2)}\\
&\deleted{=1\unsure{-}2\nu_2+4[\unsure{\nu_2}-2\nu_2^2+\nu_3(1}\\
&=1+4\nu_3-6\nu_2-24\nu_2\nu_3
\end{aligned}
\]

\deleted{We \unsure{could} take}
\[
\deleted{(\unsure{85'|85})\qquad\mu'_0=1+4\nu_3-6\nu_2}
\]

\margin{55}
\margin{55\\\texttt{https://grothendieck.umontpellier.fr}\\Université de Montpellier, UMR CNRS 5149, Case courrier 051, Place Eugène Bataillon\\34095 Montpellier cedex 5 - France}

\scanpage{57}
\mspage{676}

If we \unsure{have} ((84) and (85))
\[
\underline{\mu}=\underline{\mu}_0+12[\unsure{\deleted{\lambda_0}\inserted{\nu_0}|\nu_0}]
=1+4\nu_3-6\nu_2+12(\nu_0-2\nu_2\nu_3)
\]
we conclude from this (by \unsure{considering} the congruences modulo $4$)
\[
\tag{86}
\mu=1-2\nu_2+4\underbrace{[\nu_3-\nu_2\unsure{-|+}3\nu_2\nu_3+3\nu_0]}_{\nu'}
\]
and \unsure{by|for} \unsure{congruence|congruences} modulo $6$
\[
\tag{87}
\mu=1-2\nu_3+6\underbrace{[\nu_3-\nu_2-2\nu_2\nu_3+2\nu_0]}_{\nu''},
\]
where the underline in $\underline{\mu}$ denotes the section in $\mathbb O_{12}^*$, the non-underlined \inserted{$\mu$} \deleted{\unsure{a}} \inserted{the} \inserted{\unsure{corresponding|corresponding}} section of $\mathbf G_m=\mathbb O_1^*$. We will indeed have
\[
\begin{aligned}
2\nu'-\nu_2&=2\nu_3-3\nu_2-6\nu_2\nu_3+6\nu_0\qquad\unsure{''}\\
3\nu''-\nu_3&=2\nu_3-3\nu_2-6\nu_2\nu_3+6\nu_0
\end{aligned}
\]

In summary

\uline{Proposition} Let us introduce on \inserted{the smooth \unsure{group|group} scheme} $\mathbb O_{12}^*$ the functions \inserted{$\underline{\mu}\mapsto$} $\unsure{\nu_2|\xi_2},\nu_3\unsure{,}$ \deleted{$\unsure{\xi_3|\varepsilon_3}$ \unsure{and}} (with values in $\{0,1\}_{S_0}\subset\inserted{\mathbf G_a=}\mathbb O_1$), $\nu',\nu''$ \deleted{$\unsure{\mu_3|\nu_3}$} (with values in $\unsure{\mathbb O_1|\mathbf G_a}$) and $\mu$ (with values in $\mathbf G_m\inserted{^{\unsure{*}}=\mathbb O_1^*\subset\mathbb O_1}$) \unsure{defined} by the conditions
\[
\tag{88}
\left\{\begin{array}{ll}
\mu={}&\text{canonical \unsure{image} of }\underline{\mu}\text{ in }\mathbb O_1^*\text{ under}\\
&\text{the homomorphism of Groups }\mathbb O_N^*\unsure{\longrightarrow|\mathrel{-}}\mathbb O_1^*\\[4pt]
\begin{gathered}
\inserted{(-1)^{\nu_2}}\\[-2pt]
\deleted{\text{\unsure{the image}}}\\[-2pt]
\deleted{(\unsure{\varepsilon_2|\xi_2})}
\end{gathered}={}&\text{canonical \unsure{image} of }\underline{\mu}\text{ in }(\mathbb Z/4\mathbb Z)^*_{S_0}\\[4pt]
(-1)^{\nu_3}={}&\rule{5cm}{0.4pt}\quad(\mathbb Z/6\mathbb Z)^*_{S_0}
\end{array}\right.
\]
whence
\[
\tag{89}
\deleted{\unsure{\Bigg\{}}\quad
\left\{\begin{aligned}
\underline{\mu}&=\deleted{\unsure{\mu_0}}(-5)^{\nu_2}(+5)^{\nu_3}+12[\unsure{\nu_0|\lambda_0}]\\
&=(1+4\nu_3)(1-6\nu_2)+12\nu_0\\
&=1+4\nu_3-6\nu_2+12\unsure{[}\nu_0-2\nu_2\nu_3]
\end{aligned}\right.
\]
for a well-determined section $\nu_0$ of $\mathbf G_a=\mathbb O_1$

\margin{56}
\margin{56\\\texttt{https://grothendieck.umontpellier.fr}\\Université de Montpellier, UMR CNRS 5149, Case courrier 051, Place Eugène Bataillon\\34095 Montpellier cedex 5 - France}
